\chapter{The Heat Equation}

\section{Heat flow in a rod}

A thin rod of length $L$ has temperature $u(x,t)$. The heat equation is
\[ u_t = k\,u_{xx}, \qquad 0<x<L,\ t>0, \]
where $k$ is the thermal diffusivity. The initial temperature is $u(x,0)=f(x)$. The BCs decide which series we use.

\begin{center}
\renewcommand{\arraystretch}{1.4}
\begin{tabular}{@{}lll@{}}
\toprule
Ends & BC & Series in $x$ \\
\midrule
Held at $0^\circ$ & $u(0,t)=u(L,t)=0$ & sine series \\
Insulated & $u_x(0,t)=u_x(L,t)=0$ & cosine series \\
One held, one insulated & e.g.\ $u(0,t)=0$, $u_x(L,t)=0$ & odd-harmonic sine series \\
Held at other temperatures & $u(0,t)=T_1$, $u(L,t)=T_2$ & subtract the steady state first \\
\bottomrule
\end{tabular}
\end{center}

\section{Ends at zero temperature}

\begin{methodbox}{Heat equation with $u(0,t)=u(L,t)=0$}
Separation gives $X''+\lambda X=0$ and $T'+k\lambda T=0$. The BCs give $\lambda_n = (n\pi/L)^2$ and $X_n = \sin\frac{n\pi x}{L}$. The time equation gives $T_n = e^{-k(n\pi/L)^2t}$. So
\[ u(x,t) = \sum_{n=1}^\infty B_n\,e^{-k(n\pi/L)^2t}\,\sin\frac{n\pi x}{L},\qquad B_n = \frac2L\int_0^Lf(x)\sin\frac{n\pi x}{L}\,\dd x. \]
\end{methodbox}

\textbf{Meaning.} Unlike the wave equation, nothing oscillates in time. Each mode decays exponentially, and \emph{higher modes decay much faster} (rate $\propto n^2$). After a short time only the first mode survives. The temperature tends to zero as $t\to\infty$.

\begin{examplebox}{A hot rod is cooled at the ends}
A rod of length $L=1$ m starts at $100^\circ$C everywhere. At $t=0$ both ends are put at $0^\circ$C. The diffusivity is $k=0.01$ m$^2$/s. Find $u(x,t)$.

$B_n = 2\displaystyle\int_0^1 100\sin n\pi x\,\dd x = \dfrac{200(1-\cos n\pi)}{n\pi}$, which is $\dfrac{400}{n\pi}$ for odd $n$ and $0$ for even $n$. So
\[ u(x,t) = \frac{400}{\pi}\sum_{n\ \text{odd}}\frac{1}{n}\,e^{-0.01\,n^2\pi^2t}\,\sin n\pi x. \]
\textbf{Centre temperature at $t=10$ s.} Only $n=1$ matters: $\dfrac{400}{\pi}e^{-0.0987\cdot 10}\sin\dfrac\pi2 = 127.3\times 0.3727 \approx 47.4^\circ$C. (The $n=3$ term is about $-0.006$.)
\end{examplebox}

\begin{notebox}
The decay time of the first mode is $\tau = \dfrac{L^2}{k\pi^2}$. For this rod, $\tau = 1/(0.01\pi^2) \approx 10.1$ s. Doubling the length makes the cooling \emph{four} times slower.
\end{notebox}

\section{Insulated ends}

\begin{methodbox}{Heat equation with $u_x(0,t)=u_x(L,t)=0$}
The eigenfunctions are $\cos\frac{n\pi x}{L}$, $n = 0,1,2,\dots$ (including the constant for $n=0$). So
\[ u(x,t) = \frac{a_0}{2} + \sum_{n=1}^\infty a_n\,e^{-k(n\pi/L)^2t}\cos\frac{n\pi x}{L},\qquad a_n = \frac2L\int_0^Lf(x)\cos\frac{n\pi x}{L}\,\dd x. \]
The constant term $a_0/2$ is the \textbf{average} of $f$. It does not decay. As $t\to\infty$ the whole rod reaches this average temperature, because no heat can leave.
\end{methodbox}

\begin{examplebox}{Insulated rod}
$f(x)=x$ on $(0,L)$ with insulated ends. $a_0 = \dfrac2L\displaystyle\int_0^Lx\,\dd x = L$, so the final temperature is $L/2$. For $n\ge1$: $a_n = \dfrac2L\int_0^Lx\cos\dfrac{n\pi x}{L}\,\dd x = \dfrac{2L\big((-1)^n-1\big)}{n^2\pi^2}$, which is $-\dfrac{4L}{n^2\pi^2}$ for odd $n$.
\[ u(x,t) = \frac L2 - \frac{4L}{\pi^2}\sum_{n\ \text{odd}}\frac{1}{n^2}\,e^{-k(n\pi/L)^2t}\cos\frac{n\pi x}{L}. \]
\end{examplebox}

\section{Ends at non-zero temperatures}

\begin{methodbox}{Steady-state subtraction}
For $u(0,t)=T_1$ and $u(L,t)=T_2$:
\begin{enumerate}
  \item Find the \textbf{steady state} $u_s(x)$: it satisfies $u_s''=0$ with the same BCs, so $u_s = T_1 + (T_2-T_1)\dfrac xL$.
  \item Write $u = u_s + v$. Then $v$ solves the heat equation with $v(0,t)=v(L,t)=0$ and $v(x,0) = f(x) - u_s(x)$.
  \item Solve for $v$ with the sine series, and add $u_s$ back.
\end{enumerate}
\end{methodbox}

\begin{derivbox}{Why subtracting the steady state works}
Let $u_s(x)$ satisfy $u_s'' = 0$ with the boundary values $u_s(0)=T_1$ and $u_s(L)=T_2$. Then $u_s = T_1 + (T_2-T_1)x/L$. Define $v = u - u_s$. Since $u_s$ does not depend on $t$,
\[ v_t = u_t = ku_{xx} = k\big(u_s'' + v_{xx}\big) = kv_{xx}, \]
and $v(0,t) = T_1 - T_1 = 0$, $v(L,t) = T_2 - T_2 = 0$. So $v$ solves the heat equation with \emph{zero} end values, which is the problem we already know how to solve. The initial condition is $v(x,0) = f(x) - u_s(x)$.
\end{derivbox}

\begin{examplebox}{One end heated}
$L=1$, initial temperature $0$, $u(0,t)=0$ and $u(1,t)=100$. Then $u_s = 100x$ and $v(x,0) = -100x$.

$B_n = 2\displaystyle\int_0^1(-100x)\sin n\pi x\,\dd x$. Since $\displaystyle\int_0^1x\sin n\pi x\,\dd x = \dfrac{(-1)^{n+1}}{n\pi}$, we get $B_n = \dfrac{200(-1)^n}{n\pi}$.
\[ u(x,t) = 100x + \frac{200}{\pi}\sum_{n=1}^\infty\frac{(-1)^n}{n}\,e^{-kn^2\pi^2t}\sin n\pi x. \]
As $t\to\infty$ the temperature becomes the straight line $100x$.
\end{examplebox}

\section{Infinite rod}

For a rod with no ends, the Fourier transform (Chapter 8) turns $u_t = ku_{xx}$ into $U_t = -kw^2U$, so $U = F(w)e^{-kw^2t}$. Inverting gives

\begin{defbox}{Solution on $-\infty<x<\infty$}
\[ u(x,t) = \frac{1}{\sqrt{4\pi kt}}\int_{-\infty}^{\infty}f(\xi)\,e^{-(x-\xi)^2/(4kt)}\,\dd\xi. \]
Each point of the initial temperature spreads out like a Gaussian bell whose width grows like $\sqrt{kt}$.
\end{defbox}

\begin{derivbox}{The solution on an infinite rod}
Take the Fourier transform in $x$ of $u_t = ku_{xx}$. Let $U(w,t) = \int u(x,t)e^{-iwx}\,\dd x$. By the derivative rule (Chapter 8), $u_{xx}\to -w^2U$, and $u_t\to U_t$. So
\[ U_t = -kw^2U \ \Longrightarrow\ U(w,t) = F(w)\,e^{-kw^2t}, \]
where $F(w)$ is the transform of $f(x) = u(x,0)$. To invert, we need the function $g(x,t)$ whose transform is $e^{-kw^2t}$. Take $g = \dfrac{1}{\sqrt{4\pi kt}}e^{-x^2/(4kt)}$ and substitute $x = \sqrt{4kt}\,s$:
\[ \int g\,e^{-iwx}\,\dd x = \frac1{\sqrt\pi}\int e^{-s^2}e^{-iw\sqrt{4kt}\,s}\,\dd s = e^{-(w\sqrt{4kt})^2/4} = e^{-kw^2t}, \]
using the Gaussian transform from Chapter 8. A product of transforms is the transform of a convolution, so $u = f*g$:
\[ u(x,t) = \frac{1}{\sqrt{4\pi kt}}\int_{-\infty}^{\infty}f(\xi)\,e^{-(x-\xi)^2/(4kt)}\,\dd\xi. \]
\end{derivbox}

\begin{practicebox}
\begin{enumerate}
  \item $L=\pi$, $k=1$, ends at $0$, $u(x,0) = \sin x + 2\sin 3x$. Find $u(x,t)$.
  \item $L=\pi$, $k=1$, insulated ends, $u(x,0) = 1 + \cos 2x$. Find $u(x,t)$ and the final temperature.
  \item $L=1$, ends at $0$, $u(x,0) = x(1-x)$. Find $B_n$.
  \item A rod of length 2 m has ends held at $20^\circ$C and $80^\circ$C. Find the steady state.
  \item Find the time for the first mode to fall to $1/e$ of its starting value for $L=1$ m and $k=0.01$ m$^2$/s.
\end{enumerate}
\end{practicebox}

\begin{answerbox}
(1) $u = e^{-t}\sin x + 2e^{-9t}\sin 3x$. \quad (2) $u = 1 + e^{-4t}\cos 2x$; the final temperature is $1$. \quad (3) $B_n = 2\displaystyle\int_0^1x(1-x)\sin n\pi x\,\dd x = \dfrac{8}{n^3\pi^3}$ for odd $n$, and $0$ for even $n$. \quad (4) $u_s = 20 + 30x$. \quad (5) $\tau = L^2/(k\pi^2) \approx 10.1$ s.
\end{answerbox}
